% --- Archon physics formalization source begin ---
% archon:physics
% archon:covers PhyXMiniProblems/problem_phyx_mini_0442.lean
% archon:source-report reports/phyx_mini/problem_phyx_mini_0442.source.json

\chapter{Physics problem phyx\_mini\_0442}
\label{ch:PhyXMiniProblems_problem_phyx_mini_0442}

\paragraph{Problem source.}
\#\# Physical scenario

An object is placed against the center of a spherical mirror and then moved 70 cm from it along the central axis as the image distance \$i\$ is measured. Figure gives \$i\$ versus object distance \$p\$ out to \$p\_s = 40 \textbackslash{}text\{cm\}\$.

\#\# Question

What is \$i\$ for \$p = 70 \textbackslash{}text\{cm\}\$?

\#\# Answer choices

- A: A: 22 cm

- B: B: 24 cm

- C: C: 26 cm

- D: D: 28 cm

\#\# Figure caption (auxiliary; use the image as primary evidence)

The image contains a graph with the following characteristics:

- **Axes**: 
  - The x-axis is labeled \textbackslash{}( p \textbackslash{}, (\textbackslash{}text\{cm\}) \textbackslash{}).
  - The y-axis is labeled \textbackslash{}( i \textbackslash{}, (\textbackslash{}text\{cm\}) \textbackslash{}).

- **Plot**:
  - There is a dark curve on the graph that starts off with positive values and approaches zero as it moves from left to right, then sharply increases or decreases past the \textbackslash{}( x \textbackslash{})-axis around \textbackslash{}( p = 0 \textbackslash{}).

- **Labels**:
  - The x-axis has two labels: "0" near the intersection with the y-axis and another label \textbackslash{}( p\_s \textbackslash{}) further out along the axis.
  - The y-axis also has the label "0" indicating where it intersects the x-axis.
  
- **Grid**:
  - The graph features a grid with both vertical and horizontal lines for reference.

The curve is symmetrical around the y-axis but sharply diverges near the  \textbackslash{}( p = 0 \textbackslash{}) region.

\#\# Dataset metadata

Geometrical Optics; Physical Model Grounding Reasoning, Multi-Formula Reasoning

\paragraph{Recorded answer/context.}
D: D: 28 cm

\paragraph{Figure/image path.}
/root/proposal\_for\_physic/science-mango/phyx\_mini\_run/phyx\_data/test\_image/442.png

\paragraph{Formalization target.}
create a compiling Lean file with sorry bodies at `PhyXMiniProblems/problem\_phyx\_mini\_0442.lean`. The Lean declarations must preserve the physical quantities, dimensions or dimensional roles, figure labels, governing-law hypotheses, and final relation expressed by this problem.
Use Mathlib/Physlib names found through LeanExplore where available. If a physics API is missing, introduce faithful local abstractions rather than scalar placeholder aliases.

\begin{theorem}[Physics formalization target]
\label{thm:physics:phyx_mini_0442:target}
\lean{PhyXMiniProblems.ProblemPhyXMini0442.problem_phyx_mini_0442}
\uses{def:physics:phyx-mini-0442:phyxminiproblems-problemphyxmini0442-lengthincentimeters, def:physics:phyx-mini-0442:phyxminiproblems-problemphyxmini0442-sphericalmirrorexperiment, def:physics:phyx-mini-0442:phyxminiproblems-problemphyxmini0442-matchesproblemandintendedgraph, def:physics:phyx-mini-0442:phyxminiproblems-problemphyxmini0442-hasphysicalmirrordistances, def:physics:phyx-mini-0442:phyxminiproblems-problemphyxmini0442-graphasymptoterepresentsfocaldistance, def:physics:phyx-mini-0442:phyxminiproblems-problemphyxmini0442-satisfiesgaussiansphericalmirrorequation, def:physics:phyx-mini-0442:phyxminiproblems-problemphyxmini0442-matchesanswerchoice}
The assigned autoformalize agent should translate this physics problem into Lean declarations in the covered file, with theorem and lemma proof bodies written as `by sorry`.
\end{theorem}
\begin{proof}
This is an autoformalization task, not a proof task. Produce statements that can later be proved by the physics prover without weakening the physical model.
\end{proof}

% --- Archon declaration topology begin ---
\section{Declaration topology}
\begin{definition}[Optical Length]
  \label{def:physics:phyx-mini-0442:phyxminiproblems-problemphyxmini0442-opticallength}
  \lean{PhyXMiniProblems.ProblemPhyXMini0442.OpticalLength}
  A signed physical optical length, independent of the unit used to read it
\end{definition}
\begin{proof}
  This is the declared model definition of Optical Length.
\end{proof}
\begin{definition}[centimeter Unit Choices]
  \label{def:physics:phyx-mini-0442:phyxminiproblems-problemphyxmini0442-centimeterunitchoices}
  \lean{PhyXMiniProblems.ProblemPhyXMini0442.centimeterUnitChoices}
  SI base units with centimeters selected as the length unit
\end{definition}
\begin{proof}
  This is the declared model definition of centimeter Unit Choices.
\end{proof}
\begin{definition}[length In Centimeters]
  \label{def:physics:phyx-mini-0442:phyxminiproblems-problemphyxmini0442-lengthincentimeters}
  \lean{PhyXMiniProblems.ProblemPhyXMini0442.lengthInCentimeters}
  \uses{def:physics:phyx-mini-0442:phyxminiproblems-problemphyxmini0442-opticallength, def:physics:phyx-mini-0442:phyxminiproblems-problemphyxmini0442-centimeterunitchoices}
  The signed scalar readout of an optical length in centimeters
\end{definition}
\begin{proof}
  This is the declared model definition of length In Centimeters.
\end{proof}
\begin{definition}[Spherical Mirror Kind]
  \label{def:physics:phyx-mini-0442:phyxminiproblems-problemphyxmini0442-sphericalmirrorkind}
  \lean{PhyXMiniProblems.ProblemPhyXMini0442.SphericalMirrorKind}
  The optical type of a spherical mirror, viewed from the incident-light side
\end{definition}
\begin{proof}
  This is the declared model definition of Spherical Mirror Kind.
\end{proof}
\begin{definition}[Principal Axis Orientation]
  \label{def:physics:phyx-mini-0442:phyxminiproblems-problemphyxmini0442-principalaxisorientation}
  \lean{PhyXMiniProblems.ProblemPhyXMini0442.PrincipalAxisOrientation}
  Orientation of the central optical axis in the intended graph and experiment
\end{definition}
\begin{proof}
  This is the declared model definition of Principal Axis Orientation.
\end{proof}
\begin{definition}[Graph Axis]
  \label{def:physics:phyx-mini-0442:phyxminiproblems-problemphyxmini0442-graphaxis}
  \lean{PhyXMiniProblems.ProblemPhyXMini0442.GraphAxis}
  The two coordinate axes of the mirror-distance graph
\end{definition}
\begin{proof}
  This is the declared model definition of Graph Axis.
\end{proof}
\begin{definition}[Distance Graph Quantity]
  \label{def:physics:phyx-mini-0442:phyxminiproblems-problemphyxmini0442-distancegraphquantity}
  \lean{PhyXMiniProblems.ProblemPhyXMini0442.DistanceGraphQuantity}
  Physical quantities attached to the graph labels p and i
\end{definition}
\begin{proof}
  This is the declared model definition of Distance Graph Quantity.
\end{proof}
\begin{definition}[Mirror Distance Graph]
  \label{def:physics:phyx-mini-0442:phyxminiproblems-problemphyxmini0442-mirrordistancegraph}
  \lean{PhyXMiniProblems.ProblemPhyXMini0442.MirrorDistanceGraph}
  \uses{def:physics:phyx-mini-0442:phyxminiproblems-problemphyxmini0442-opticallength, def:physics:phyx-mini-0442:phyxminiproblems-problemphyxmini0442-graphaxis, def:physics:phyx-mini-0442:phyxminiproblems-problemphyxmini0442-distancegraphquantity}
  The intended i-versus-p graph. pSubS is the dimensionful quantity printed as p\_s; verticalAsymptoteObjectDistance records the object-distance coordinate at which the plotted curve diverges. The image-distance response is stored as a physical length for every object distance, including distances beyond the displayed range
\end{definition}
\begin{proof}
  This is the declared model definition of Mirror Distance Graph.
\end{proof}
\begin{definition}[Spherical Mirror Experiment]
  \label{def:physics:phyx-mini-0442:phyxminiproblems-problemphyxmini0442-sphericalmirrorexperiment}
  \lean{PhyXMiniProblems.ProblemPhyXMini0442.SphericalMirrorExperiment}
  \uses{def:physics:phyx-mini-0442:phyxminiproblems-problemphyxmini0442-opticallength, def:physics:phyx-mini-0442:phyxminiproblems-problemphyxmini0442-sphericalmirrorkind, def:physics:phyx-mini-0442:phyxminiproblems-problemphyxmini0442-principalaxisorientation, def:physics:phyx-mini-0442:phyxminiproblems-problemphyxmini0442-mirrordistancegraph}
  The spherical mirror, its intended graph, and the two axial positions needed by the question. axialDisplacement is the object's physical motion from its initial position; requestedObjectDistance is the resulting value of p. No value is assigned here to the requested image distance
\end{definition}
\begin{proof}
  This is the declared model definition of Spherical Mirror Experiment.
\end{proof}
\begin{definition}[Matches Problem And Intended Graph]
  \label{def:physics:phyx-mini-0442:phyxminiproblems-problemphyxmini0442-matchesproblemandintendedgraph}
  \lean{PhyXMiniProblems.ProblemPhyXMini0442.MatchesProblemAndIntendedGraph}
  \uses{def:physics:phyx-mini-0442:phyxminiproblems-problemphyxmini0442-lengthincentimeters, def:physics:phyx-mini-0442:phyxminiproblems-problemphyxmini0442-sphericalmirrorexperiment}
  Problem-statement data and the explicit intended-graph calibration. The object begins at p = 0, moves 70 cm, and the graph's horizontal scale mark is p\_s = 40 cm. Because the supplied raster is unrelated, the condition that the vertical asymptote bisects the displayed scale is a transparent conditional input for the intended graph, not a claimed raster readout. No value of i at p = 70 cm occurs here
\end{definition}
\begin{proof}
  This is the declared model definition of Matches Problem And Intended Graph.
\end{proof}
\begin{definition}[Has Physical Mirror Distances]
  \label{def:physics:phyx-mini-0442:phyxminiproblems-problemphyxmini0442-hasphysicalmirrordistances}
  \lean{PhyXMiniProblems.ProblemPhyXMini0442.HasPhysicalMirrorDistances}
  \uses{def:physics:phyx-mini-0442:phyxminiproblems-problemphyxmini0442-lengthincentimeters, def:physics:phyx-mini-0442:phyxminiproblems-problemphyxmini0442-sphericalmirrorexperiment}
  Positivity conditions selecting the physical concave-mirror branch
\end{definition}
\begin{proof}
  This is the declared model definition of Has Physical Mirror Distances.
\end{proof}
\begin{definition}[Graph Asymptote Represents Focal Distance]
  \label{def:physics:phyx-mini-0442:phyxminiproblems-problemphyxmini0442-graphasymptoterepresentsfocaldistance}
  \lean{PhyXMiniProblems.ProblemPhyXMini0442.GraphAsymptoteRepresentsFocalDistance}
  \uses{def:physics:phyx-mini-0442:phyxminiproblems-problemphyxmini0442-sphericalmirrorexperiment}
  For an ideal paraxial spherical mirror, the vertical asymptote of the i-versus-p graph occurs at p = f. This connects a generic graph feature to the mirror's unknown focal length and contains no requested image distance
\end{definition}
\begin{proof}
  This is the declared model definition of Graph Asymptote Represents Focal Distance.
\end{proof}
\begin{definition}[Satisfies Gaussian Spherical Mirror Equation]
  \label{def:physics:phyx-mini-0442:phyxminiproblems-problemphyxmini0442-satisfiesgaussiansphericalmirrorequation}
  \lean{PhyXMiniProblems.ProblemPhyXMini0442.SatisfiesGaussianSphericalMirrorEquation}
  \uses{def:physics:phyx-mini-0442:phyxminiproblems-problemphyxmini0442-opticallength, def:physics:phyx-mini-0442:phyxminiproblems-problemphyxmini0442-sphericalmirrorexperiment}
  The signed Gaussian spherical-mirror equation 1 / f = 1 / p + 1 / i is written, away from p = f, in the division-free homogeneous form f * (p + i) = p * i. Requiring it for every unit choice makes the relation independent of the centimeter readouts used by the graph
\end{definition}
\begin{proof}
  This is the declared model definition of Satisfies Gaussian Spherical Mirror Equation.
\end{proof}
\begin{definition}[Answer Choice]
  \label{def:physics:phyx-mini-0442:phyxminiproblems-problemphyxmini0442-answerchoice}
  \lean{PhyXMiniProblems.ProblemPhyXMini0442.AnswerChoice}
  Labels of the four displayed answer choices
\end{definition}
\begin{proof}
  This is the declared model definition of Answer Choice.
\end{proof}
\begin{definition}[image Distance In Centimeters]
  \label{def:physics:phyx-mini-0442:phyxminiproblems-problemphyxmini0442-answerchoice-imagedistanceincentimeters}
  \lean{PhyXMiniProblems.ProblemPhyXMini0442.AnswerChoice.imageDistanceInCentimeters}
  \uses{def:physics:phyx-mini-0442:phyxminiproblems-problemphyxmini0442-answerchoice}
  The image-distance readout printed beside each answer label, in centimeters
\end{definition}
\begin{proof}
  This is the declared model definition of image Distance In Centimeters.
\end{proof}
\begin{definition}[Matches Answer Choice]
  \label{def:physics:phyx-mini-0442:phyxminiproblems-problemphyxmini0442-matchesanswerchoice}
  \lean{PhyXMiniProblems.ProblemPhyXMini0442.MatchesAnswerChoice}
  \uses{def:physics:phyx-mini-0442:phyxminiproblems-problemphyxmini0442-opticallength, def:physics:phyx-mini-0442:phyxminiproblems-problemphyxmini0442-lengthincentimeters, def:physics:phyx-mini-0442:phyxminiproblems-problemphyxmini0442-answerchoice}
  Exact agreement of a physical image distance with a displayed choice
\end{definition}
\begin{proof}
  This is the declared model definition of Matches Answer Choice.
\end{proof}
\begin{lemma}[focal Length In Centimeters eq twenty]
  \label{lem:physics:phyx-mini-0442:phyxminiproblems-problemphyxmini0442-focallengthincentimeters-eq-twenty}
  \lean{PhyXMiniProblems.ProblemPhyXMini0442.focalLengthInCentimeters_eq_twenty}
  \uses{def:physics:phyx-mini-0442:phyxminiproblems-problemphyxmini0442-lengthincentimeters, def:physics:phyx-mini-0442:phyxminiproblems-problemphyxmini0442-sphericalmirrorexperiment, def:physics:phyx-mini-0442:phyxminiproblems-problemphyxmini0442-matchesproblemandintendedgraph, def:physics:phyx-mini-0442:phyxminiproblems-problemphyxmini0442-hasphysicalmirrordistances, def:physics:phyx-mini-0442:phyxminiproblems-problemphyxmini0442-graphasymptoterepresentsfocaldistance}
  Under the explicit intended-graph calibration, p\_s = 40 cm and the focal asymptote lies at half that scale, so the mirror has focal length 20 cm
\end{lemma}
\begin{proof}
  The conclusion follows from the governing relations and figure readouts named in the statement of focal Length In Centimeters eq twenty.
\end{proof}
% --- Archon declaration topology end ---
% --- Archon physics formalization source end ---
